\section{Finite-memory optimality at the least feasible budget}\label{sec:boundary-memory}
Boundary attainment in Theorem~\ref{thm:boundary-attainment} concerns the compact space of full-history plans. It does not assert that the optimum belongs to any finite-label class. We resolve that distinction on a class of prepared heralded channel experiments. The result is a necessary-and-sufficient memory criterion, an explicit finite-dimensional quantum obstruction, and quantitative approximation at the same boundary budget. It is not an additional assumption of the general compactness theorem.

\subsection{Prepared heralded channels and a boundary rigidity identity}
Let $0<p<1$, put $q=1-p$, and fix a nonempty finite failure alphabet $Y$ with probabilities $\pi_y>0$ summing to one. There is one globally legal unit-cost command. Model $\theta$ has a prepared density matrix $\rho_\theta$ on a fixed finite-dimensional space and channels $\Phi_{\theta,y}$. Its report instrument is
\begin{equation}\label{eq:herald-instrument}
 \mathcal I_{\theta,s}(A)=pA,\qquad
 \mathcal I_{\theta,y}(A)=q\pi_y\Phi_{\theta,y}(A)\quad(y\in Y).
\end{equation}
The herald $s$ is productive; a failure $y$ forbids commitment. Rejection of $s$ leaves the physical state unchanged. The observer cannot inspect the protected physical state. After commitment to $r$ in one fixed nonempty finite dictionary $C$, the future audit has conditional loss observable $0\preccurlyeq F_{\theta,r}\preccurlyeq LI$. Any separately charged fixed final-audit cost is added to all budgets below. No unpriced channel, phase input or clock is supplied to the observer.

For $w=y_1\cdots y_n\in Y^*$, including the empty word, set
\[
 \rho_\theta(w)=\Phi_{\theta,y_n}\cdots\Phi_{\theta,y_1}(\rho_\theta),\quad
 \lambda(w)=pq^{|w|}\prod_{i=1}^{|w|}\pi_{y_i},\quad
 \ell_\theta(w,r)=\tr(F_{\theta,r}\rho_\theta(w)).
\]
Every word has positive mass, $\sum_w\lambda(w)=1$, and $\sum_{|w|\ge T}\lambda(w)=q^T$. Probabilities of visible failures are independent of both the hidden model and the private controller coins. The channels themselves may change the prepared state noncommutatively.

\begin{lemma}[The least budget forces first-herald commitment]\label{lem:boundary-herald}
The least feasible precommit expected-call budget is $N_{\min}=1/p$. For every proper finite-cost policy, let $J$ count productive reports which it rejects. Then, in every model,
\begin{equation}\label{eq:herald-cost-identity}
 p\E_\theta\tau=1+\E_\theta J.
\end{equation}
Consequently a policy is feasible at $N_{\min}$ if and only if it commits at the first herald, up to null histories. At that budget the only choice is its decision distribution $b_w\in\Delta(C)$ after each failure word $w$, and its risk is
\begin{equation}\label{eq:boundary-risk-functional}
 \mathcal R_\theta(b)=\sum_w\lambda(w)\sum_{r\in C}b_w(r)\ell_\theta(w,r).
\end{equation}
\end{lemma}
\begin{proof}
The event that call $i$ is made is predictable, and the conditional herald probability at that call is $p$. Summing the nonnegative indicators gives $\E_\theta\sum_{i\le\tau}\one_{\{s_i=s\}}=p\E_\theta\tau$. Exactly one of these heralds is accepted; all others are rejected. This proves the identity without an optional-stopping integrability assumption. It implies the lower $1/p$, attained by first-herald commitment. Equality forces the nonnegative integer $J$ to vanish almost surely. Conversely that rule has a geometric call count of mean $1/p$. Conditional on a failure word, the physical state is the displayed channel product and is independent of the private tape. The independent, unobserved audit is performed only after the decision. This gives the risk formula.
\end{proof}

\subsection{A robust decision quotient and its exact memory content}
Assume the model space $\Theta$ is compact metric and the prepared states, channels and audit observables are continuous in $\theta$. A singleton is permitted; an arbitrary prior is not part of the data. Write $V=\min_b\sup_\theta\mathcal R_\theta(b)$ at $N_{\min}$. For a probability measure $\mu$ on $\Theta$ put
\[
 \ell_\mu(w,r)=\int\ell_\theta(w,r)\,d\mu(\theta),\quad
 a_\mu(w)=\min_{r\in C}\ell_\mu(w,r),\quad
 V(\mu)=\sum_w\lambda(w)a_\mu(w).
\]

\begin{theorem}[Boundary minimax and the regularity criterion]\label{thm:boundary-regularity}
The primal optimum and a least-favorable measure are attained, and
\begin{equation}\label{eq:boundary-minimax}
 V=\max_{\mu\in\mathcal P(\Theta)}V(\mu).
\end{equation}
For an attaining $\mu_*$, every robust optimum has zero expected decision slack:
\begin{equation}\label{eq:boundary-slack}
 \int\mathcal R_\theta(b)\,d\mu_*(\theta)-V
 =\sum_w\lambda(w)\sum_r b_w(r)
          [\ell_{\mu_*}(w,r)-a_{\mu_*}(w)]\ge0.
\end{equation}
Suppose the minimizer $a(w)\in C$ of $\ell_{\mu_*}(w,\cdot)$ is unique for every word. Then $a$ is the unique robust optimal decision rule, and a finite stochastic autonomous observer attains $V$ if and only if the relation
\begin{equation}\label{eq:decision-congruence}
 w\equiv_a v\quad\Longleftrightarrow\quad
          a(wz)=a(vz)\text{ for every }z\in Y^*
\end{equation}
has finite index. The least number of live labels attaining $V$ is exactly that index. Random transitions and randomized readouts cannot reduce this exact count.

At each prescribed finite live-label bound $M$, the best boundary risk is attained. If the index is infinite, that best risk is strictly larger than $V$ for every finite $M$, although finite-label policies converge to $V$ at the same budget. A prefix implementation uses at most
$1+\sum_{j=0}^{T-1}|Y|^j$ live labels and has excess risk at most $Lq^T$. Its output dictionary remains $C$. For any implementation with $M$ live labels and at most $m$ actual decoder values, event stopping uses at most $\max(M,m)$ total labels and state stopping uses at most $M+m$, by Lemma~\ref{lem:roles}.
\end{theorem}
\begin{proof}
The product $\mathcal B=\prod_{w\in Y^*}\Delta(C)$ is compact and convex. The tail of every risk series is at most $Lq^T$, uniformly in $b$ and $\theta$. Each truncated expression is jointly continuous and affine in the appropriate variable. The full payoff on $\mathcal B\times\mathcal P(\Theta)$ is therefore continuous and bilinear, with both factors compact convex. The elementary compact minimax theorem applies \cite{Sion58}. For completeness, it can also be obtained here from the finite-prefix games: truncate at $T$, separate the convex hull of the finite-dimensional model loss vectors from a strict lower orthant, and take compact subsequential limits of priors and plans. The uniform error $Lq^T$ passes both inequalities to the limit. Minimizing for a fixed measure chooses a minimizing decision separately at every word and gives $V(\mu)$. This proves equality and both attainments. Nonnegative summation proves the slack identity. At a robust optimum its left side is at most zero and hence is zero. Since every $\lambda(w)>0$, uniqueness forces $b_w$ to be the point mass at $a(w)$. Existence of a robust optimum then proves that this Bayes rule is itself robustly optimal; a general Bayes minimizer need not have this property when ties occur.

At the boundary an $M$-label stochastic observer is represented by one stochastic transition matrix $Q_y$ for each failure letter and one decoder distribution $d_s\in\Delta(C)$ at each live label. Its initial live label is fixed. Unreachable rows can be completed arbitrarily. Zero slack says that every label reached with positive probability after $w$ decodes $a(w)$ with probability one. Choose one positive-probability successor at each pair $(s,y)$ and the uniquely prescribed output on each reached label. This gives a deterministic observer with no more labels and the same output at every word: its path always had positive probability in the stochastic observer. Thus exact random realization implies deterministic realization.

If a deterministic machine sends $w,v$ to the same label, every common continuation sends them through identical transitions and gives the same output. Hence they are equivalent in \eqref{eq:decision-congruence}. Any exact machine has at least its number of equivalence classes. Conversely this relation is a right congruence: $w\equiv_a v$ implies $wy\equiv_a vy$. Its classes, with transition $[w]\mapsto[wy]$ and decoder $a(w)$, form an exact deterministic observer whenever the index is finite. This is the standard residual-function construction for finite automata, applied to the optimal decision rule rather than to the full predictive state.

For fixed $M$, all $Q_y$ and $d_s$ lie in a finite product of compact simplexes. Each finite risk sum is continuous in their entries. The same uniform geometric tail makes full risks, and their supremum over $\Theta$, continuous. A minimum is attained. It can equal $V$ only if an exact machine exists, proving the strict gap for infinite index. Finally follow an optimal full-history rule on words of length less than $T$, storing that finite tree, and use one fallback decoder label thereafter. The policy still stops at the first herald and costs exactly $1/p$. Only the terminal mass $q^T$ can change loss, proving the approximation and state bound. This argument uses no interior budget slack.
\end{proof}

The quotient in \eqref{eq:decision-congruence} is specific to this boundary task and this fixed dictionary. It is not a common predictive quotient preserving every experiment or future audit. The minimax argument is classical; the new content is its combination with exact minimum-cost commitment, zero-slack determinization, and a necessary-and-sufficient autonomous memory count. Finite-index residual characterizations and unary periodicity are established automata theory \cite{Mereghetti24}; their use does not establish a new priority claim for those tools.

\begin{proposition}[Ties for one preparation, and why robust ties differ]\label{prop:boundary-ties}
For a single prepared model, without uniqueness, an exact finite stochastic observer exists if and only if there is a finite-index right congruence $\sim$ on $Y^*$ such that
\[
 \bigcap_{w\in A}\arg\min_{r\in C}\ell(w,r)\ne\varnothing
             \quad\text{for every equivalence class }A.
\]
The least number of live labels is the least number of classes among such congruences. For several unknown models the same zero-slack condition is necessary for an optimal stochastic observer under a least-favorable measure, but deterministic selection from tied Bayes actions need not preserve robust optimality.
\end{proposition}
\begin{proof}
For one preparation, every wordwise minimizing rule is globally optimal. Zero slack makes the decoder support of any reached stochastic label a subset of every optimal-action set at its reaching words. Choose one such output and a positive successor for each failure letter. Determinization preserves optimality. The partition of words by its reached deterministic label is a right congruence with the displayed nonempty intersections. Conversely choose an output in each intersection and run the congruence-class observer. Minimizing the number of classes proves the count.

For the distinction, take two models with identical channels and reports and losses $(0,1)$ and $(1,0)$ on a two-point dictionary, independent of the word. The robust value is $1/2$, attained with one live label and a fair terminal decision. Under the equally weighted least-favorable measure both actions are tied at loss $1/2$. Either constant deterministic selection has worst loss one. Thus Bayes zero slack alone cannot license robust determinization in the presence of ties.
\end{proof}

\begin{proposition}[A finite distinguishing-depth lower certificate]\label{prop:boundary-depth-gap}
In the unique-action case of Theorem~\ref{thm:boundary-regularity}, suppose the index is infinite and $m=|C|\ge2$. Fix $M$. There is a finite $H_M$ such that every deterministic $M$-label observer disagrees with $a$ at a word of length at most $H_M$. Put
\[
 \gamma_H=\min_{|w|\le H,\ r\ne a(w)}
       [\ell_{\mu_*}(w,r)-\ell_{\mu_*}(w,a(w))]>0,
 \qquad \pi_{\min}=\min_y\pi_y.
\]
Every stochastic $M$-label boundary observer obeys
\begin{equation}\label{eq:boundary-quantitative-gap}
 \sup_\theta\mathcal R_\theta-V
       \ge\frac{p\gamma_{H_M}}{m}
                    \left(\frac{q\pi_{\min}}{M}\right)^{H_M}>0.
\end{equation}
This is a finite distinguishing certificate, not a sharp asymptotic formula for the optimal intermediate-memory risk.
\end{proposition}
\begin{proof}
There are finitely many deterministic transition and output tables on $M$ named labels. Each table fails at some finite word because the residual index is infinite. The maximum of their first distinguishing depths is finite. From any stochastic table choose a largest-probability successor for each $(s,y)$, of probability at least $1/M$, and a largest-probability output at each label, of probability at least $1/m$. The resulting deterministic table fails at a word $w$ of length $n\le H_M$. In the stochastic table the selected internal path and output occur with probability at least $M^{-n}/m$ conditional on this word. Their slack is at least $\gamma_{H_M}$. The word's terminal mass is at least $p(q\pi_{\min})^n$. Keep just that term in \eqref{eq:boundary-slack}; since $q\pi_{\min}/M<1$, lowering it to the exponent $H_M$ gives the bound. Worst-model risk is no smaller than the least-favorable average.
\end{proof}

\subsection{An algebraic qubit with no finite-memory boundary optimum}
A quantum orbit supplies a finite-dimensional physical realization, not an arbitrary nonregular string placed in an external clock. Take one failure letter and a single prepared qubit,
\begin{equation}\label{eq:qubit-boundary-data}
 \rho_0=|0\rangle\langle0|,\qquad
 U=\frac1{\sqrt5}\begin{pmatrix}2&-1\\1&2\end{pmatrix},\qquad
 \Phi(A)=UAU^*,\qquad C=\{-1,1\}.
\end{equation}
On commitment the future audit measures $Z=\operatorname{diag}(1,-1)$, with outcomes $z\in\{-1,1\}$, and scores $(z-r)^2$. Thus $F_r=2I-2rZ$ and $L=4$. Write $\vartheta$ for the Bloch rotation angle, so $\cos\vartheta=3/5$, $\sin\vartheta=4/5$, and put $c_n=\cos(n\vartheta)$. In particular $c_0=1$, $c_1=3/5$, and $c_{n+1}=(6/5)c_n-c_{n-1}$.

\begin{theorem}[Boundary nonattainment by every finite classical memory]\label{thm:qubit-boundary}
For every $0<p<1$, the least budget is $1/p$ and the full-history risk there is
\begin{equation}\label{eq:qubit-boundary-value}
 V_p=2-2p\sum_{n=0}^\infty q^n|c_n|.
\end{equation}
Its unique optimal rule is $a_n=\operatorname{sgn}c_n$. It has infinite residual index. Hence no finite stochastic autonomous observer attains $V_p$ at this budget. This remains true with arbitrary real transition probabilities and randomized terminal decisions.

Let $V_{p,M}$ be the optimal risk with $M$ live labels and the same two-point dictionary. Each such minimum exists, and
\begin{equation}\label{eq:qubit-memory-gap}
 0<V_{p,M}-V_p\le4q^M.
\end{equation}
The exact one-live-label optimum is
\begin{equation}\label{eq:qubit-one-label}
 V_{p,1}=2-\frac{2p(1-3q/5)}{1-6q/5+q^2}.
\end{equation}
At $p=1/2$ it is $12/13$, while the full-history value is the strictly smaller convergent series in \eqref{eq:qubit-boundary-value}. All specified primitive entries are algebraic. The bound \eqref{eq:boundary-quantitative-gap} is an explicit positive lower with $m=2$, $\pi_{\min}=1$, and $\gamma_H=4\min_{0\le n\le H}|c_n|$.
\end{theorem}
\begin{proof}
After $n$ failures the prepared state is $U^n\rho_0U^{*n}$, whose audit mean is $c_n$. The expected squared loss of decision $r$ is $2-2rc_n$. Lemma~\ref{lem:boundary-herald} forces first-herald commitment, so independent pointwise minimization gives \eqref{eq:qubit-boundary-value}.

The number $e^{i\vartheta}$ is not a root of unity. Otherwise $2\cos\vartheta=6/5$ would be a rational algebraic integer, hence an integer, a contradiction. Thus $\vartheta/\pi$ is irrational; no $c_n$ is zero, and each residue subsequence $(n_0+jr)\vartheta$ is dense modulo $2\pi$ for every positive integer $r$. The latter elementary rotation fact follows, for example, by the pigeonhole approximation of a nonzero arbitrarily small rotation and its successive multiples. Each residue subsequence therefore visits both open semicircles on which cosine has opposite signs. The optimal sign sequence cannot be ultimately periodic. A finite deterministic unary observer always follows a finite initial path and then a cycle, so it produces an ultimately periodic decision sequence. Theorem~\ref{thm:boundary-regularity} now excludes every finite stochastic observer and proves the strict gap and its attainment at fixed $M$.

Use $M$ live labels to store $n=0,\ldots,M-1$, incrementing on failure until the last label, which thereafter repeats. Decode the correct sign at these $M$ labels and commit at the first herald. The rule is exact for $n<M$ and its excess loss is bounded by four on the tail of mass $q^M$. Its cost is still exactly $1/p$. This proves the upper in \eqref{eq:qubit-memory-gap}.

For one live label the expected output is one constant $b\in[-1,1]$. Its risk is $2-2bA$, where
\[
 A=p\sum_{n\ge0}q^nc_n
   =p\operatorname{Re}\frac1{1-qe^{i\vartheta}}
   =\frac{p(1-3q/5)}{1-6q/5+q^2}>0.
\]
Thus $b=1$ is optimal and gives \eqref{eq:qubit-one-label}. Substituting $p=q=1/2$ gives $12/13$. The failure and herald maps are genuine completely positive trace-nonincreasing maps with trace-preserving sum. They use one declared command and a protected terminal audit; the proof never reads that audit early or turns the qubit into observer-owned memory.
\end{proof}

\begin{corollary}[A compact uncertain-preparation family]\label{cor:polarization-boundary}
Replace the known initial qubit by
$\rho_\theta=(I+\theta Z)/2$, with unknown fixed $\theta\in[\eta,1]$ and $0<\eta\le1$. Keep the same instrument, dictionary and audit. The robust boundary value is
\[
 V_{p,[\eta,1]}=2-2\eta p\sum_{n\ge0}q^n|c_n|.
\]
A least-favorable preparation is $\rho_\eta$, and the unique robust optimal rule is again $\operatorname{sgn}c_n$. It is unattainable by every finite stochastic autonomous memory, while the same-budget prefix approximation remains valid. Thus the obstruction is not tied to a chosen Bayesian prior or to one pure initial preparation. It concerns this declared preparation family, not all qubit density matrices.
\end{corollary}
\begin{proof}
At preparation $\theta$ the conditional mean is $\theta c_n$. The sign rule minimizes loss word by word for every positive $\theta$, and its total risk decreases with $\theta$. Its worst value is therefore the displayed one at $\eta$. Every common policy is lower-bounded by its risk in that same preparation, whose unique minimizer is the sign rule. This proves the robust value, the least-favorable measure $\delta_\eta$, and uniqueness without an assumed prior. Apply Theorem~\ref{thm:boundary-regularity}. If $\eta<1$, the family includes both mixed and pure states; no inference is made at $\eta=0$, where the least-favorable audit is a fair bit and uniqueness fails.
\end{proof}

\begin{corollary}[Exact classification on unitary orbits]\label{cor:unitary-period}
For any finite-dimensional unitary failure channel, fixed prepared state, fixed finite dictionary and audit observables, suppose the minimizing decision along the orbit is unique at every nonnegative integer time. A finite-label boundary optimum exists if and only if the optimal decision sequence is periodic from time zero. Its least live-label count is the least positive period.
\end{corollary}
\begin{proof}
The residual theorem first gives equivalence with ultimate periodicity. If the sequence has eventual period $r$, there are integers $k_j\to\infty$ with $U^{rk_j}\to I$ in operator norm, by compactness of the finite-dimensional unitary torus. For each fixed $n$, the states at times $n+rk_j$ and $n+r+rk_j$ converge to those at $n$ and $n+r$. Uniqueness among finitely many decisions gives a strict gap at each of these two limiting states, so their optimal decisions eventually agree with the limiting decisions. Eventual $r$-periodicity makes these decisions equal. Therefore $a_n=a_{n+r}$ for every $n\ge0$. A cycle of length the least period realizes the rule, and every finite unary machine has an eventual cycle at least that long.
\end{proof}

\subsection{Dissipation and the strict decision-gap condition}
Finite memory is not confined to periodic unitary orbits. A stable decision at the limiting state gives a different positive realization mechanism.

\begin{proposition}[Uniform contraction with a strict limiting decision gap]\label{prop:contractive-boundary}
Consider a single prepared heralded model with finitely many failure letters. Suppose there is a density matrix $\rho_\infty$ and constants $C_0<\infty$, $0<\alpha<1$ such that
\[
 \|\rho(w)-\rho_\infty\|_1\le C_0\alpha^{|w|}
                \quad(w\in Y^*).
\]
Suppose one dictionary value $r_*$ has a strict limiting gap
\[
 g=\min_{r\ne r_*}\tr[(F_r-F_{r_*})\rho_\infty]>0.
\]
For any nonnegative integer $T$ with $2LC_0\alpha^T<g$, there is an exact optimal boundary observer with at most
$1+\sum_{n=0}^{T-1}|Y|^n$ live labels, where the empty sum is zero. All words of length at least $T$ use the same decision $r_*$. Ties at earlier words are harmless in this single-model statement. In particular, dissipative channels with a certified contraction and a strict limiting task gap can have an exact finite-memory optimum at the least feasible cost.
\end{proposition}
\begin{proof}
For every $r\ne r_*$ and every word of length $n$, trace duality gives
\[
 \ell(w,r)-\ell(w,r_*)
 \ge g-\|F_r-F_{r_*}\|_\infty
                    \|\rho(w)-\rho_\infty\|_1
 \ge g-2LC_0\alpha^n.
\]
This is strictly positive at $n\ge T$. Store all shorter words in a finite tree and merge their depth-$T$ continuations into one label decoding $r_*$. Choose a least-loss value separately at each shorter word. Every decision minimizes its actual conditional loss, and the rule stops at the first herald, so it attains the single-model optimum at cost $1/p$. The displayed number counts the entire prefix tree and the common tail label. No limiting state is supplied to the observer during execution; it is an offline certificate of its finite program.
\end{proof}

The limiting gap is a real restriction, not a consequence of mixing. The following family keeps a strictly positive contraction rate but retains the infinite exact-memory obstruction.
\begin{corollary}[Noisy rotating qubits: contraction is not sufficient]\label{cor:noisy-boundary}
In Theorem~\ref{thm:qubit-boundary}, replace the failure channel by
\[
 \Phi_\alpha(A)=\alpha UAU^*+(1-\alpha)\tr(A)I/2,
                    \qquad 0\le\alpha\le1.
\]
For $0<\alpha\le1$, the exact boundary value and one-live-label value are
\begin{align*}
 V_{p,\alpha}&=2-2p\sum_{n\ge0}(q\alpha)^n|c_n|,\\
 V_{p,\alpha,1}&=2-\frac{2p(1-3q\alpha/5)}{1-6q\alpha/5+q^2\alpha^2}.
\end{align*}
No finite stochastic observer attains $V_{p,\alpha}$, while
\[
 0<V_{p,\alpha,M}-V_{p,\alpha}
       \le \frac{4p(q\alpha)^M}{1-q\alpha}.
\]
For $0<\alpha<1$ the failure channel is uniformly trace-norm contractive to $I/2$. At $\alpha=0$, however, one live label is exactly optimal, with value $2-2p$. Thus a positive mixing margin alone cannot replace the strict limiting decision gap in Proposition~\ref{prop:contractive-boundary}.
\end{corollary}
\begin{proof}
Depolarization commutes with unitary conjugation, so
$\rho_n=I/2+\alpha^n U^n(\rho_0-I/2)U^{*n}$ and its audit mean is $\alpha^nc_n$. The unique minimizing decision is therefore the same nonperiodic sign sequence whenever $\alpha>0$. Apply Theorem~\ref{thm:boundary-regularity}. A wrong decision at time $n$ has excess loss at most $4\alpha^n$; summing the tail from $M$ gives the displayed bound. For a constant decision, replace $q$ by $q\alpha$ in the elementary cosine series, retaining the factor $p$, to obtain the one-label value. Trace differences of states have zero trace and contract by exactly the factor $\alpha$ under this channel. At $\alpha=0$ every positive-length failure word has state $I/2$ and both decisions are tied, while at the empty word the decision $+1$ is optimal. The constant $+1$ rule attains $2-2p$. At $I/2$ the two limiting losses agree, explaining why the preceding positive-gap theorem does not apply.
\end{proof}

\subsection{The exact two-label frontier for a noisy quantum orbit}
The qubit application also admits a complete nontrivial finite-memory optimization. The next formula allows every rotation angle and damping strength, rather than fixing the algebraic example's angle.

\begin{theorem}[Exact two-live-label quantum prediction law]\label{thm:two-label-quantum}
Consider the same prepared qubit, herald and squared Pauli audit, with a Bloch rotation angle $\vartheta$, $c=\cos\vartheta\in[-1,1]$, and depolarizing strength $0\le\alpha\le1$. Put $x=q\alpha$ and
\[
 f_c(t)=\frac{1-ct}{1-2ct+t^2},\qquad A=p f_c(x).
\]
The optimal one-live-label risk is $2-2A$. The exact optimum among all stochastic observers with at most two live labels is $2-2S_2$, where the following cases determine $S_2$ and an attaining transition matrix. On a herald the first label decodes $+1$ and the second $-1$, except in the constant-decision cases.

If $x=0$, then $S_2=p$ and one label suffices. For $x>0$ and $0\le c\le1$, let $s=\sqrt{1-c^2}$ and $t_*=c/(1+s)$. Then
\[
 S_2=\begin{cases}
 A,&x\le t_*,\\
 p(1+s)/s-A,&x>t_*.
 \end{cases}
\]
The second case never occurs at $c=1$. In the first case use a constant $+1$ decision. In the second case an optimal failure matrix is
\[
 Q=\begin{pmatrix}t_*/x&1-t_*/x\\0&1\end{pmatrix}.
\]
For $x>0$ and $-1\le c<0$, write $k=-c$, $s=\sqrt{1-k^2}$, and
\[
 t=\min\left\{x,\max\left\{0,
       \frac{x+k-s\sqrt{1+2kx+x^2}}{1+kx}\right\}\right\}.
\]
Then
\[
 S_2=A+\frac{2[p f_c(-t)-A]}{1+t/x},\qquad
 Q=\begin{pmatrix}0&1\\t/x&1-t/x\end{pmatrix}
\]
is attaining. All formulas include ties and degenerate transitions. The first compulsory call is included in the cost, which remains exactly $1/p$ in every case.

In particular, for the irrational algebraic orbit with $c=3/5$, $p=1/2$ and $\alpha=1$, the exact two-live-label risk is $43/52$, attained by
$Q=\left(\begin{smallmatrix}2/3&1/3\\0&1\end{smallmatrix}\right)$.
For this same angle, two live labels improve on one if and only if $q\alpha>1/3$. This is a phase boundary for the declared random-horizon prediction task, not the unresolved intermediate-memory deterministic-release law.
\end{theorem}
\begin{proof}
At the least budget, Lemma~\ref{lem:boundary-herald} reduces every feasible observer to first-herald commitment. Its two-state failure matrix has the form
$Q=\left(\begin{smallmatrix}1-a&a\\b&1-b\end{smallmatrix}\right)$, $0\le a,b\le1$, from initial label zero. Write $u_i\in[-1,1]$ for the expected decoded sign at label $i$. Its risk is
\[
 2-2v\cdot u,\qquad
 v=p\sum_{n\ge0}x^n\cos(n\vartheta)e_0Q^n.
\]
The series is absolutely convergent because $x<1$. If $a+b=0$, then $v=(A,0)$. Otherwise put $\lambda=1-a-b$ and $z=a/(a+b)$. The elementary two-state decomposition gives
\[
 v_0=(1-z)A+zB,\quad v_1=z(A-B),\qquad
 B=p f_c(x\lambda)>0.
\]
Also $A>0$, so $v_0>0$. Optimizing the two decoder signs therefore gives reward $A$ if $B\le A$, and $A+2z(B-A)$ otherwise. For a fixed $\lambda$, the largest feasible $z$ is one when $\lambda\ge0$ and $1/(1-\lambda)$ when $\lambda<0$.

For $c\ge0$, differentiation gives
$f_c'(t)=(c-2t+ct^2)/(1-2ct+t^2)^2$. Its maximum on $[-x,x]$ occurs at $\min\{x,t_*\}$. When $x\le t_*$, $B\le A$ for every matrix. Otherwise the maximizing $\lambda=t_*/x$ is nonnegative and permits $z=1$. Thus the upper $2\max B-A$ is attainable, and $f_c(t_*)=(1+s)/(2s)$ gives the formula.

For $c=-k<0$, nonnegative $\lambda$ is optimized at zero. For negative $\lambda$ set $t=-x\lambda\in[0,x]$ and choose $z=1/(1+t/x)$. The improvement divided by two is
\[
 \frac{p x}{1+2kx+x^2}
       \frac{x+k-(1+kx)t}{1-2kt+t^2}.
\]
Its derivative has the sign of
$(1+kx)t^2-2(x+k)t+2k(x+k)-(1+kx)$.
The larger root exceeds $x$ because $(x+k)/(1+kx)>x$ for $k>0$. Consequently the maximum on $[0,x]$ is at the smaller root clipped to this interval; that root is the displayed $t$. The endpoint zero includes the best nonnegative-$\lambda$ choice. This proves global optimality over all two-state matrices and all randomized decoders. Substitution gives $A=7/13$, $\max B=9/16$ in the specified algebraic case, hence $2-2(2\cdot9/16-7/13)=43/52$. The strict improvement condition follows from the positive-angle cases. No numerical optimization is used in this proof.
\end{proof}

\subsection{A finite-precision synthesis bound for unary rational channels}
There is also a positive computational conclusion at the boundary, rather than only a nonattainment theorem. This statement concerns one known prepared channel and a fixed finite dictionary; the full robust synthesis problem is not asserted to have the same complexity.

\begin{theorem}[Polynomial prefix synthesis at a bounded herald rate]\label{thm:boundary-algorithm}
Suppose $Y$ has one letter, the prepared density matrix and the channel superoperator in real Hermitian coordinates have rational entries, and the $m$ loss observables have rational entries and lie in $[0,LI]$. Let $p$ and a requested $0<\epsilon<L$ be rational. From these data one can construct an autonomous observer of exact cost $1/p$ and risk at most $V+\epsilon$, using
\[
 T=\max\left\{1,\left\lceil\frac{\log(L/\epsilon)}{-\log(1-p)}\right\rceil\right\}
\]
live labels. The integer $T$ may instead be chosen by exact successive rational tests of $L(1-p)^T\le\epsilon$, without evaluating a logarithm.

After a common-denominator preprocessing, suppose the data entries have denominator dividing $D$ and numerators and $D$ have at most $B$ bits; include the rational herald rate, loss bound and tolerance in the input size. The synthesis uses $O(T(d^4+md^2))$ rational coordinate operations on integers with $O(TB+\log d)$ bits, in addition to the scalar stopping test. A direct implementation has $O(d^2(TB+\log d))$ bits of working storage, apart from its output program, and a transition/readout table of $O(T(\log(T+1)+\log(m+1)))$ bits. Thus, for a herald rate bounded away from zero, this is a polynomial bit-time construction in the rational input size and $\log(1/\epsilon)$, using standard integer arithmetic. It is not a polynomial bound uniform as $p\downarrow0$, a matching lower bound, or a complexity claim for continuum-model robust optimization.
\end{theorem}
\begin{proof}
Compute $\rho_n=\Phi^n(\rho_0)$ for $0\le n<T$. At each step evaluate the $m$ rational numbers $\tr(F_r\rho_n)$ and choose a least one, resolving ties by the fixed dictionary order. Use a chain of $T$ live labels; on failure advance until the last label, which repeats. At a herald commit the decision stored at that label. The rule is optimal on all terminal words of length less than $T$. Its excess loss is at most $L$ on the remaining terminal mass $q^T$, while its call count is still geometric with mean $1/p$. This proves the guarantees.

A channel application is a matrix--vector product in $d^2$ real Hermitian coordinates, requiring $O(d^4)$ scalar operations. Evaluating all decisions costs $O(md^2)$ more. Use real and imaginary matrix-entry coordinates, with off-diagonal entries paired, so all physical state coordinates have absolute value at most one. After $n$ applications their denominator divides $D^{n+1}$. Keeping this common denominator explicitly gives integer numerators of $O((n+1)B)$ bits; a loss inner product adds only $O(B+\log d)$ bits. The same size bound follows directly during multiplication and summation, and comparison of decisions uses integer comparisons at the shared denominator. Store two state vectors and the current best decision; the stated working bound includes the growing state numerators. The input itself is read-only and not included in working storage. The output table stores a successor and a dictionary index per live label, with no stored floating phase or external clock. Standard integer addition and multiplication have polynomial bit cost, giving the final claim. If the exact scalar test is used, the denominator and numerator of $q^T$ also have $O(TB)$ bits. The displayed bound on $T$ is logarithmic in precision only when $p$ has the stated positive lower bound.
\end{proof}

The qubit in Theorem~\ref{thm:qubit-boundary} belongs to this rational-superoperator class: although $U$ has algebraic entries, its conjugation map has rational coefficients with denominator five. Its exact optimum requires unbounded classical memory, but its boundary-optimal risk has efficiently computable finite-memory approximants. The stored number of live-label bits is $\lceil\log_2T\rceil$; this does not replace the separate program and arithmetic storage accounts above.

The algebraic rotating-qubit idea and nonperiodic quantum acceptance events have established antecedents in quantum automata \cite{Mereghetti24}. We do not claim priority for that mechanism. Here it is compiled into a paid, single-command prediction experiment whose least feasible cost is attained, whose full-history bounded-risk optimum is attained, and whose exact risk is nevertheless unattainable with every fixed finite classical memory. The same boundary admits exponentially accurate finite-label approximations. The transition-table account excludes the fixed numerical dictionary and primitive description, whose encoding lengths must also be added when measuring total program size. These facts distinguish optimum attainment, fixed-memory gaps, and density, rather than inferring one from another. They do not settle the sharp intermediate-memory phase-type timing curve of Theorem~\ref{thm:timer-frontier}, the general boundary-density question, or a universal unknown-model predictive quotient.
